\lesson{2}{Sep 09 2021 Thu (08:20:00)}{Measuring Matter}{Unit 1}

\subsubsection*{Measurements and Units}

All measurements are made of \bf{2} parts:

\begin{itemize}
  \item A number.
  \item A unit.
\end{itemize}

\begin{definition}[Matter]
  Anything that takes up space, whether big or small, is \bf{Matter}.
\end{definition}

\begin{definition}[Volume]
  In science, the amount of space an object occupies is called \bf{volume}
\end{definition}

\begin{definition}[Mass]
  \bf{Mass} is a measure of the amount of matter in an object. This means, the more \bf{Matter} an object has, the object's \bf{Mass} increases. \\
  Scientists measure the \bf{Mass} of an object using the base \bf{Unit Gram}.
\end{definition}

\subsubsection*{Systems of Measurements}

The \bf{U.S. Customary System (English System)} consists of measurements we use every day to describe:

\begin{itemize}
  \item How much we have.
  \item What size we want.
  \item How far we want to go.
  \item Etc \ldots
\end{itemize}

The \bf{United States} is the one of three countries in the \bf{whole world} that use this system.

\begin{center}
  \begin{table}[h]
    \begin{tabular}{ |r|c|l| } 
      \hline
      Type & U.S. Customary System & Metric System \\
      \hline
      \multirow{3}{4em}{} Length      & inch, foot, mile 		& meter 			\\ 
      Mass/Weight & ounce, pound, ton 	& gram 			    \\ 
      Volume 		& pint, quart, gallon 	& liter 		    \\ 
      Temperature	& degrees Fahrenheit 	& degrees Celsius   \\ 
      \hline
    \end{tabular}

    \caption{The \bf{U.S. Customary System (English System)}}
    \label{table:english-system}
  \end{table}
\end{center}

The \bf{Metric System} is a system of measuring units based on the power of
$10$. It's the preferred system used in science measurement because so many
countries already use metrics as their standard measurement system.

One reason for the widespread use of the \bf{Metric System} is that it
seems to simplify the number of measurements. Instead of multiple units for the
same measurements, such as:

\begin{itemize}
  \item Inches.
  \item Feet.
  \item Miles.
\end{itemize}

for length, the \bf{Metric System} uses only one word for the length, the
\bf{Meter}.

For larger or smaller measurements, metric prefixes are used
rather than changing the unit of measure or using fractions, like quarter of an
inch. A prefix works with any unit of measure in the \bf{Metric System}:

\begin{itemize}
  \item Meter.
  \item Gram.
  \item Liter.
\end{itemize}

\begin{center}
  \begin{table}[h]
    \begin{tabular}{ |r|c|c|l| } 
      \hline
      Prefix & Symbol & Multiplier & Example \\
      \hline
      \multirow{3}{4em}{} Kilo-   & k     & $1,000$   & Kilometer (km) 	\\ 
      Hecto- 	& h     & $100$ 	& Hectometer (hm)   \\ 
      Deca- 	& da 	& $10$ 		& Decameter (dam)   \\ 
      Unit	&  	    & $1$ 		& Meter (m) 		\\ 
      Deci-	& d 	& $0.1$ 	& Decimeter (dm)    \\ 
      Centi-	& c 	& $0.01$ 	& Centimeter (cm)   \\ 
      Milli-	& m 	& $0.001$   & Millimeter (mm)   \\ 
      \hline
    \end{tabular}

    \caption{The \bf{Metric System Prefix References}}
    \label{table:metric-system}
  \end{table}
\end{center}

The \bf{International System of Measurement of Units} is the primary system
used by scientists and engineers. It will also be used in \bf{Chemistry}.
In this system, scientists have specified what units should be used as the
seven base units of measurement from which all other units are derived.

\begin{center}
  \begin{table}[h]
    \begin{tabular}{ |r|l| } 
      \hline
      Quantity (Symbol) & Name of Unit (Symbol) \\
      \hline
      \multirow{3}{4em}{} Length (l)                      & 	Meter (m) 		\\ 
      Mass (m)                        & 	Kilogram (kg)   \\
      Time (t) 					    & 	Second (s) 		\\
      Electric Current (i) 			& 	Ampere (A) 		\\
      Thermodynamic Temperature (T) 	& 	Kelvin (k) 		\\
      Amount of Substance (n) 		& 	Mole (mol) 		\\
      Luminous Intensity ($l_v$) 		& 	Candela (cd) 	\\
      \hline
    \end{tabular}

    \caption{The \bf{Metric System Prefix References}}
    \label{table:metric-system}
  \end{table}
\end{center}

\subsubsection*{Converte between or within Measurement Systems}

\begin{definition}[Dimensional Analysis] Analyzing the relationships between
  different quantities, identify their units of measurement, and convert their
  quantities for equal comparison.

  \bf{Dimensional Analysis} uses a ratio, or a conversion factor to change
  units of measurement while maintaining the value of measurement.

  In a \bf{Conversion Factor}, the \bf{numerator} and \bf{denominator} equal the same amount, just in different units.

  For example, to convert $3$ miles into feet, you would multiply $3$ miles by
  a conversion factor created from the units of miles and feet. Because $1$
  mile equals $5,280$ feet, two conversion factors can be created

  \begin{align}
    1 = \frac{1 \rm{ mile}}{5,280 \rm{ feet}} \rm{ or } 1 = \frac{5,280 \rm{ feet}}{1 \rm{ mile}}
  \end{align}

  Multiplying any number, such as $3$ miles, by either of these fractions does
  not change the number's overall value.
\end{definition}

\subsubsection*{Using the conversion factor for converting derived units}

\begin{marginfigure}
  \begin{tikzpicture}[every edge quotes/.append style={auto, text=blue}]
    \pgfmathsetmacro{\cubex}{3}
    \pgfmathsetmacro{\cubey}{3}
    \pgfmathsetmacro{\cubez}{3}
    \draw [draw=blue, every edge/.append style={draw=blue, densely dashed, opacity=.5}, fill=magenta]
      (0,0,0) coordinate (o) -- ++(-\cubex,0,0) coordinate (a) -- ++(0,-\cubey,0) coordinate (b) edge coordinate [pos=1] (g) ++(0,0,-\cubez)  -- ++(\cubex,0,0) coordinate (c) -- cycle
      (o) -- ++(0,0,-\cubez) coordinate (d) -- ++(0,-\cubey,0) coordinate (e) edge (g) -- (c) -- cycle
      (o) -- (a) -- ++(0,0,-\cubez) coordinate (f) edge (g) -- (d) -- cycle;
    \path [every edge/.append style={draw=blue, |-|}]
      (b) +(0,-5pt) coordinate (b1) edge ["3cm"'] (b1 -| c)
      (b) +(-5pt,0) coordinate (b2) edge ["3cm"] (b2 |- a)
      (c) +(3.5pt,-3.5pt) coordinate (c2) edge ["3cm"'] ([xshift=3.5pt,yshift=-3.5pt]e)
      ;
  \end{tikzpicture}
  \caption{Cube that's $3 \times 3 \times 3$}
  \label{fig:3-3-3-cube}
\end{marginfigure}

\begin{definition}[Derived Units]
  Some measurements you are familiar with, such as speed, contain more than one \bf{Unit}.

  A very common \bf{Derived Unit} is \bf{Volume}. For liquids and gasses,
  \bf{Volume} tends to be measured in liters. But, for regular-shaped solids,
  measurements and calculations can determine their \bf{Volume}. To calculate
  the \bf{Volume} of the cube in Figure~\ref{fig:3-3-3-cube}, just raise the
  side measurement to the third power.

  \begin{align}
    V &= (3 cm)^3 \\
      &= (3 \times 3 \times 3) cm \\
      &= 27 cm^3 
  \end{align}
\end{definition}

Now, let's learn how to convert \bf{Derived Units}:

\newpage

\begin{center}
  \begin{table}[h]
    \begin{tabular}{ |r|c|l| } 
      \hline
      Length & Weight & Liquid Capacity \\
      \hline
      \multirow{3}{4em}{} 1 foot = 12 inches      & 1 pound = 16 ounces 	& 1 cup = 8 fluid ounces    \\
      1 yard = 3 feet 		& 1 ton = 2,000 pounds 	& 1 pint = 2 cups           \\
      1 mile = 5,280 feet 	& 						& 1 quart = 2 pints         \\
      1 mile = 1,760 yards	& 						& 1 gallon = 4 quarts       \\
      \hline
    \end{tabular}

    \caption{U.S. Customary Units of Measure}
    \label{table:U.S.-customary-units-of-measure}
  \end{table}
\end{center}

\begin{center}
  \begin{table}[h]
    \begin{tabular}{ |r|c|l| } 
      \hline
      Length & Weight/Mass & Liquid Capacity \\
      \hline
      \multirow{3}{4em}{} 1 foot = 2.54 centimeters   & 1 pound = 0.454 kilograms & 1 gallon = 3.785 liters           \\
      1 meter = 39.37 inches 	    & 1 ton = 2.2 pounds 		& 1 liter = 0.264 gallons  	 	    \\ 
      1 mile = 1.609 kilometers 	& 1 ounce = 28.35 grams		& 1 liter = 1,000 cubic centimeters \\ 
      1 kilometer = 0.6214 miles	& 							& 1 milliliter = 1 cubic centimeter \\ 
      \hline
    \end{tabular}

    \caption{Metric Units of Measure}
    \label{table:metric-units-of-measure}
  \end{table}
\end{center}

\begin{center}
  \begin{table}[h]
    \begin{tabular}{ |r|c|c|l| } 
      \hline
      Prefix & Symbol & Multiplier & Example \\
      \hline
      \multirow{3}{4em}{} Kilo- 	& k     & $1,000$   & Kilometer (km) 	\\ 
      Hecto- 	& h 	& $100$ 	& Hectometer (hm) \\ 
      Deca- 	& da	& $10$ 		& Decameter (dam) \\ 
      Unit	&  		& $1$ 		& Meter (m) 			\\ 
      Deci-	& d 	& $0.1$ 	& Decimeter (dm) 	\\ 
      Centi-	& c 	& $0.01$ 	& Centimeter (cm) \\ 
      Milli-	& m 	& $0.001$   & Millimeter (mm) \\ 
      \hline
    \end{tabular}

    \caption{The \bf{Metric System Prefix References}}
    \label{table:metric-system}
  \end{table}
\end{center}

\begin{example}[Fractional Unit]
  If you start with a number that has a
  factional unit, there may be more steps to the \bf{Dimensional Analysis}.
  It doesn't have to be harder, you still want to make sure that you cancel out
  all the units except the units you want in your final answer. Look at the
  following problem to see how to change the numerator and/or denominator of a
  fractional unit.

  In this problem, you've been asked to convert from the fractional unit
  kilometers per hour to the unit centimeters per second. This means that
  you'll need to convert the numerator and denominator to get the final answer.

  When you have a given that has a fractional unit, it can be helpful to write
  the measurement as a fraction with an one in the denominator:

  \begin{align}
    \frac{0.36 km}{1 h}
  \end{align}

  This makes it more obvious that there are two numbers that need to convert:
  $0.36 km$ and $1 h$.

  Once you know what needs to be converted, you can determine those
  relationships that will help you go from kilometers to centimeters and hours
  to seconds.

  Since you need to change both the numerator and the denominator, you can
  start with the numerator first. The one hour is still in the denominator, but
  you can ignore it until you get kilometers converted to centimeters:

  \begin{align}
    \frac{0.36 \rm{ \sout{km}}}{1 h} \times
    \frac{1,000 \rm{ \sout{m}}}{1 \rm{ \sout{km}}}
    \times \frac{100 cm}{1 \rm{ \sout{m}}}
  \end{align}

  Notice that the units that need to be canceled ($km$ and $m$ ) are all
  diagonal from each other. That leaves centimeters as the only unit that does
  not cancel. It is in the numerator, where it needs to be for the final
  answer.

  Now that we've changed the numerator from the unit $km$ to $cm$ on top, you
  can continue the problem by changing the denominator from $hr$ to $sec$:

  \begin{align}
    \frac{36,000 cm}{1 \rm{ \sout{h}}} \times \frac{1 \rm{ \sout{h}}}{60 \rm{ \sout{min}}} \times \frac{1 \rm{ \sout{min}}}{60 sec}
  \end{align}

  Notice that you need to put hours on the top of the next conversion in order
  to cancel out with ours in the denominator from $hr$ to $sec$

  Notice that you need to put hours on the top of the next conversion in order
  to cancel out with hours in the denominator of the given. In the next step,
  the minutes unit is canceled to leave seconds in the denominator -- where it
  needs to be for the final answer:

  \begin{align}
    \frac{36,000 cm}{1 \rm{ \sout{h}}} \times \frac{1 \rm{ \sout{h}}}{60 \rm{ \sout{min}}} \times \frac{1 \rm{ \sout{min}}}{60 sec} &= 10 \frac{cm}{s} \\
                                                                                                                                    &= 0.36 \frac{km}{h} \\
                                                                                                                                    &= 10 \frac{cm}{s}
  \end{align}
\end{example}

\begin{example}[Powered Unit] Some units are raised to a power.
  \bf{Area} and \bf{Volume} of solids are the common ones.

  Area is measured in squared length, such as the square meter, $m^2$, while
  \bf{Volume} is measured in cubic length, such as the cubic meter, $m^3$.
  Because one cubic centimeter equals $1$ milliliter, that conversion is quite
  simple. However, there are times you must convert from one cubic meter to a
  cubic centimeter. Let's look at an example:

  Convert the \bf{Volume} $5.0 m^3$ to the unit $cm^3$.

  Just like any unit conversion, we look at the starting and ending units, then
  think about what we know that can help us convert the units. You haven't
  learned an equivalent relationship  between cubic metres and cubic
  centimeters, but you know the relationship between the meters and
  centimeters:

  \begin{align}
    1 m = 100 cm
  \end{align}

  We simply use the relationship we already know and cube the entire thing.
  This makes the fraction a three dimensional \bf{Volume} measurement
  instead of a one dimensional length measurement:

  \begin{align}
    \left( \frac{100 cm}{1 m}\right)^3
  \end{align}

  After we do that, we can setup the problem, by starting with the given
  measurement and adding the relationship we already know.

  The units $m^3$ and $m$ cannot cancel out because they are not exactly the
  same. The entire conversion factor must be cubed in order to make the units
  match:

  \begin{align}
    5.0 m^3 \times \frac{100 cm}{1 m} &= 5.0 m^3 \times \left( \frac{100 cm}{1 m} \right)^3 \\
                                      &= 5.0 m^3 \times \left( \frac{1,000,000 cm^3}{1 m^3} \right)
  \end{align}

  Now that the cubic meter units are able to cancel, it's time to calculate the
  answer. Because the entire fraction is being cubed, the numbers and units must
  all be cubed:

  \begin{align}
    5.0 m^3 \times \frac{1,000,000 cm^3}{1 m^3} \times \frac{5.0 \times 1,000,000 cm^3}{1} = 5,000,000 cm^3 \rm{ or } 5 \times 10^6 cm^3
  \end{align}

  Therefore, $5.0 m^3$ is equivalent to $5 \times 10^6 cm^3$.
\end{example}

\newpage
